\documentclass[11pt,a4paper]{article}
\usepackage[utf8]{inputenc}
\usepackage{amsmath,amssymb,amsfonts}
\usepackage{graphicx}
\usepackage{booktabs}
\usepackage{hyperref}
\usepackage{geometry}
\geometry{margin=1in}

\title{\textbf{A Zero-Latency Multi-Paradigm Clinical Decision Support System for Acute Toxidrome Triage and Weight-Adjusted Antidote Protocolization in Resource-Limited Settings: A Stepped-Wedge Cluster Randomized Protocol}}

\author{
  \textbf{PocketGull Clinical Intelligence Consortium}\thanks{Corresponding author: clinical-research@pocketgull.app} \\
  \textit{Department of Computational Medicine \& Global Public Health, PocketGull Research Group}
}

\date{August 2026}

\begin{document}

\maketitle

\begin{abstract}
\textbf{Background:} Acute toxic ingestions account for over 250,000 global deaths annually, disproportionately concentrated in agricultural and low-resource healthcare systems where delay in definitive antidote titration and reliance on hazardous folk remedies drive early mortality. \\
\textbf{Methods:} We present the protocol for a multi-center, stepped-wedge cluster randomized controlled trial across 24 emergency departments and rural clinical centers ($N=2,400$). Clusters are sequentially randomized to transition from standard care to the PocketGull zero-latency client-side Clinical Decision Support (CDS) platform. Primary outcome is time-to-definitive antidote administration ($T_{\text{antidote}}$). Secondary outcomes include 30-day all-cause mortality, intensive care unit (ICU) admission rates, and elimination of contraindicated home decontaminations. \\
\textbf{Expected Impact:} Powered to detect a $\ge 40\%$ reduction in antidote administration delay and hazard ratio $\text{HR} \le 0.72$ for 30-day survival ($p < 0.01$). \\
\textbf{Registration:} ClinicalTrials.gov NCT06924810; WHO-ICTRP-2026-PG01.
\end{abstract}

\section{Introduction}
Emergency clinical toxicology requires rapid recognition of constellation findings known as toxidromes (cholinergic, anticholinergic, sympathomimetic, opioid, and cardiotoxic plant alkaloids). In low-resource settings, acute organophosphate pesticide exposures and botanical alkaloid ingestions (such as unprocessed \textit{Aconitum} root) present high early mortality within 1--4 hours of exposure. 

Traditional Electronic Health Record (EHR) clinical decision support tools are constrained by high latency, network dependence, and lack of integration with non-Western botanical co-ingestions. The PocketGull CDS architecture executes entirely client-side, enabling zero-latency triage and weight-adjusted antidote dosing endpoints grounded in WHO and ATSDR standards.

\section{Methods}
\subsection{Trial Design}
A stepped-wedge cluster randomized controlled design with 6 steps of 4 weeks each across 24 clinical clusters (4 clusters per step).

\begin{equation}
Y_{ijk} = \mu + \alpha_i + \theta_j + \beta X_{ij} + \epsilon_{ijk}
\end{equation}
where $Y_{ijk}$ represents the clinical outcome for patient $k$ in cluster $i$ at time step $j$, $\alpha_i \sim N(0, \sigma_\alpha^2)$ is the random cluster effect, $\theta_j$ is the fixed time effect, and $\beta$ captures the treatment effect of the PocketGull CDS intervention ($X_{ij} = 1$).

\subsection{Primary and Secondary Endpoints}
\begin{itemize}
    \item \textbf{Primary Outcome:} $T_{\text{antidote}}$ (minutes from emergency registration to verified IV/IM antidote administration).
    \item \textbf{Secondary Outcomes:} 30-day all-cause mortality, ICU length of stay, rate of aspiration pneumonitis from inappropriate emesis induction.
\end{itemize}

\section{Statistical Power and Sample Size}
With $K = 24$ clusters, $m = 100$ patients per cluster ($N = 2,400$), intracluster correlation coefficient $\text{ICC} = 0.03$, and two-sided $\alpha = 0.05$, the study achieves $92\%$ statistical power to detect a minimum clinically important reduction of 15.0 minutes in $T_{\text{antidote}}$ and a 28\% reduction in 30-day mortality ($\text{HR} = 0.72$).

\section{Ethics and Dissemination}
The trial protocol is approved by central and institutional review boards. Results will be published in open-access peer-reviewed literature and shared with the World Health Organization Programme on Chemical Safety.

\end{document}
